% Original: PDF strana 24, gornji slajd.
\slajdid{02-s047}
\begin{frame}[label=02-s047]{Grupisanje proizvoda i broj čipova}
\setlength{\parskip}{0pt}
Algebarskom preraspodelom proizvoda možemo realizovati istu funkciju pomoću dvoulaznih I kola, uz povećano kašnjenje:
% element: 02-s047:e01
\[F=\bar D\bar B+D\bar C\bar A=\bar D\bar B+(D\bar C)\bar A\]
\par\vspace{3mm}
% element: 02-s047:e02
Za ovu realizaciju potrebni su:
\begin{itemize}
\item jedan čip sa invertorima;
\item jedan čip sa dvoulaznim I kolima — koristimo \textbf{tri I kola}, koja se nalaze u jednom čipu;
\item jedan čip sa dvoulaznim ILI kolima.
\end{itemize}
\par\vspace{2mm}
% element: 02-s047:e03
\Rezultat{Ukupno \textbf{3 čipa}. Troulazni proizvod realizujemo kaskadom dva dvoulazna I kola: smanjujemo broj potrebnih kućišta, \alert{\textbf{na račun povećanog kašnjenja}}.}
\end{frame}
